Everything below is fixed by the released code and its committed seeds; this appendix
collects the settings a reader would otherwise have to recover from source.

\paragraph{What ``retraining from scratch'' means here.} All fine-tuning and all
unlearning in this paper operate on LoRA adapters of rank $16$
($\alpha_{\mathrm{LoRA}} = 32$, dropout $0$, target modules \texttt{all-linear})
over a \emph{frozen} pretrained base model. ``Retrain'' therefore means: initialise a
\emph{fresh} adapter on the same frozen base and train it on
$D_{\text{keep}}$ only---the keep split with neither the organic forget data nor any
canary---using the identical optimiser, schedule and step count as the original
fine-tune. It is exact unlearning \emph{within the adapter paradigm}: the adapter that
saw the canaries is discarded and never contributes to the retrained model, and the
frozen base never saw a canary at any point. It is not retraining of the base model's
pretraining, which we never perform and do not claim. We now call this subject
``retrain (fresh adapter)'' in the tables to remove the ambiguity.

Whether adapter-level unlearning is representative of a deployment setting deserves a
direct answer, since the certificate is meant to cover the latter. The audit itself is
indifferent: it queries $M_u$ as a black box and never inspects parameters, so nothing
in Theorem~\ref{thm:null} or Theorem~\ref{thm:cert} changes if the provider unlearns by
full-parameter gradient steps instead. What the adapter setting does change is the
\emph{difficulty} of the unlearning problem the subjects face. A rank-$16$ adapter has
far fewer degrees of freedom in which to hide memorised content than a full
fine-tune, so our subjects are, if anything, being certified in a regime that favours
them; a provider doing full-parameter unlearning should expect verdicts at least as
demanding. We flag this as a limitation of the empirical study rather than of the
framework, and full-parameter runs are the natural next scaling step.

\paragraph{Models.} \texttt{gpt2} (124M), \texttt{EleutherAI/pythia-160m} (160M),
\texttt{microsoft/phi-1\_5} (1.4B), \texttt{Qwen/Qwen3-0.6B},
\texttt{Qwen/Qwen3-4B-2507}, \texttt{microsoft/Phi-4-mini-instruct},
\texttt{google/gemma-4-E2B-it}, and an in-repo 0.9M-parameter TinyGPT for the
CPU-affordable tier. Phi-1.5 LoRA fine-tuning diverges to NaN in fp16, so all Phi-1.5
runs are fp32; the remaining large models run fp16 with a shared frozen base and one
adapter per training stage.

\paragraph{Data.} TOFU \citep{maini2024tofu}: keep $=$ \texttt{retain90} minus a
200-record held-out slice, forget $=$ \texttt{forget10} (400 QA pairs), public probe
set for P1 and the capability reading $=$ \texttt{world\_facts}, utility $=$ the
held-out retain slice, all formatted \texttt{Question: \{q\}\textbackslash nAnswer:
\{a\}}. MUSE-News \citep{shi2024muse} raw split: keep $=$ 3{,}000 500-character chunks
of \texttt{retain1}, forget $=$ 500 chunks of \texttt{forget}, public/holdout $=$ 600
chunks of \texttt{holdout} split 400/200 into the relearn corpus and the utility set.

\paragraph{Canaries.} Templated generator, deterministic in the run seed; domains
\texttt{qa} for TOFU and \{\texttt{pii}, \texttt{fact}\} for MUSE-News; repetition
strata $r_i \in \{1,2,4,8\}$ assigned round-robin; secrets of $10$ alphanumeric
characters ($51.7$ bits). Cohort sizes: $384$ pairs (TOFU/GPT-2, TOFU/Pythia, all
model-axis runs), $512$ (TOFU/Phi-1.5, all MUSE-News runs, TinyGPT), $640$ (the
GPT-2 synthetic tier). Canary character share of the fine-tuning corpus is $11.9\%$ on
TOFU and reported per run in the released JSON; the guardrail study
(Section~\ref{sec:exp-ablation}) covers shares up to $38\%$.

\paragraph{Training and unlearning.} Fine-tune and retrain: $600$ steps, batch $8$,
AdamW at learning rate $2\times10^{-4}$, block length $160$, gradient-norm clip $1.0$.
Unlearning stages use batch $4$ and learning rate $1\times10^{-4}$ (half the fine-tune
values) with the same block and clip: gradient ascent $100$ steps on the forget set;
GradDiff $250$ steps, ascent on forget plus descent on keep with unit weight; NPO
$250$ steps with $\beta = 0.1$ and the fine-tuned adapter as the frozen reference;
SimNPO $250$ steps with $\beta = 0.1$ and margin $\gamma = 0$, reference-free on the
length-normalised forget NLL; ``NPO (25\% budget)'' truncates NPO to $62$ steps.
Probes: P1 relearn $=$ $100$ steps on the public corpus, which never contains a
canary; P2 $=$ symmetric per-channel round-to-nearest 4-bit quantisation of all linear
and embedding weights; P3 $=$ three adversarial prompt wrappers applied at scoring
time only. Probe budgets are fixed across all subjects and models.

\paragraph{Scoring.} $Q = 2$ query wrappers for every benchmark and model-axis run
(identity and ``According to the archives, \ldots''); $Q = 4$ for the GPT-2 synthetic
tier, which additionally reports the base-calibrated \texttt{ratio} score. Score class
$\F = \{s_{\mathrm{loss}}, s_{\mathrm{mink}}\}$ for the benchmark and model-axis runs
and $\F = \{s_{\mathrm{loss}}, s_{\mathrm{mink}}, s_{\mathrm{ratio}}\}$ for the
synthetic tier, with $k = 20\%$ for min-$k$. Scores are computed from temperature-0
forward passes and aggregated over wrappers by the mean. Verification cost is exactly
$2Qm$ forward passes per subject.

\paragraph{Two implementation notes that affect validity.} First, the stake cap and the
wealth factor must be computed from the \emph{same} clamped value of the recentred
boundary $p_0(t)$. The admissible range is $\lambda_t < c/(1-p_0(t))$, so deriving the
cap from one clamped boundary and the payoff from another lets the worst-case factor
$1 + \lambda_t(p_0(t)-1)$ fall below zero. This is reachable exactly when $m\,p_0$ is
an integer---which holds for $m = 640$ at every tolerance we report, in both
directions. An earlier build of the released code had two different clamps and did
produce non-positive raw factors on that tier; the current build shares one clamp,
carries a regression test, and we have re-run the affected tier, whose verdicts are
unchanged.

Second, Theorem~\ref{thm:cert} requires the reveal order $\pi$ to be independent of
the realised sign vectors. Every result table in this paper was produced with
\texttt{reveal\_source = "seeded"}, in which $\pi$ is seeded from the run seed---the
same seed that drives cohort generation and training---so for those runs the
independence does not strictly hold, even though nothing in the pipeline exploits the
coupling. We report them as they were computed rather than restating them under a
different reveal rule. The released code now also implements the two sources that do
satisfy the hypothesis (\texttt{vouch/verify/beacon.py}): \texttt{"beacon"} seeds $\pi$
with $H(\text{manifest root} \,\|\, r)$ for the randomness $r$ of a drand round drawn
after the commitment is published, recording the round number and value on the
certificate so a third party can recompute the order; and \texttt{"local"} draws from
operating-system entropy. A deployed audit should use the beacon, and the certificate
object now carries a \texttt{reveal} block naming the source, whether it is
third-party verifiable, and whether it meets Theorem~\ref{thm:cert}'s hypothesis.

\paragraph{Verification.} $\alpha = 0.05$; two-sided certification; sign-based
revocation for all headline tables, with the magnitude-aware arm reported separately;
finite-cohort certificate target by default and the super-population target reported
alongside wherever tolerance matters. Confidence-sequence grid $1001$ points on
$[0,1]$. Betting mixture: fixed fractions $\{0.02, 0.05, 0.1, 0.2, 0.4, 0.7\}$ of
$\lambda_{\max}$, an ONS expert with $\eta = \tfrac12$ and $A_0 = 1$, and a truncated
KT plug-in, combined with uniform prior weights. Reveal order: a permutation seeded by
the run seed, committed before any query. Tie-breaking: a committed PRNG seeded at
$20260702$, shared across every run so that tie decisions are reproducible.
Early stopping is enabled for the certificate arm; the revocation arm continues to
accumulate evidence after a crossing so that the streaming product alarm is
well defined.

\paragraph{Simulation tiers.} Validity, soundness, dependence and baselines: $2{,}000$
seeds. Power and censoring: $500$ seeds per cell with a $20{,}000$-pair horizon.
Tightness and streaming: $500$ seeds. Group-sequential boundaries use $K = 10$ equally
spaced looks with Pocock and O'Brien--Fleming spending functions, calibrated by exact
binomial dynamic programming over the partial-sum lattice at the least favourable
null.

\paragraph{Hardware and cost accounting.} The study ran on heterogeneous free and local
compute: Colab and Kaggle T4/P100 GPU sessions for the 1.4B--5.1B models, an
Apple-silicon laptop (MPS backend) for the GPT-2 and Pythia benchmark runs and the
revision-round SimNPO runs, and Colab CPU sessions plus a four-core container for the
TinyGPT tier and all simulations. Because the hardware is not matched across tiers, we
report verifier cost as a \emph{workload} in forward passes and give measured
wall-clock times only alongside the tier that produced them
(Section~\ref{sec:exp-ablation}). The comparison against retrain-based and
shadow-model evaluators is likewise stated as a ratio of workloads on the same model
and hardware---$1{,}536$ inference passes against roughly $4{,}800$ training passes per
fine-tune and $16$ fine-tunes for a shadow panel---rather than as a wall-clock
speed-up factor.
